\chapter{Gestión de errores en C-Horizon}
\label{cap:errores}

Para cada error se indicará su tipo y su posición en el archivo (fila y columna). Si el error se encuentra en una instrucción simple, se retomará la 
compilación justo donde termina la instrucción, esto es, después de ";". En caso de que el error se de en la declaración de un ámbito (una función, una 
condición, un bucle) se retomará la compilación justo al terminar ese ámbito. Es decir, se buscará el caracter "\{" y se omitirá todo el código hasta 
encontrar el "\}" que lo cierre (para ello también se llevará un contador de los "\{" sin cerrar). Algunos ejemplos de errores: 

\begin{itemize}
	\item \err{[Line 6. Column 9. Code 01] Type missmatch. Expected type \Integer.}
	\item \err{[Line 54. Column 12. Code 02] Invalid operation: divided by zero.}
\end{itemize}
%\errmsg{Code 01: Type missmatch. Expected type .} Significa que se esperaba un tipo concreto ("_") y se ha recibido otro. 
%\errmsg{Code 02: Inavlid operation} Significa que se ha realizado una operación aritmética inválida (dividir entre 0).
